\section*{AI工具使用声明}

本参赛队在竞赛过程中使用了 AI 工具，主要用于思路讨论、代码调试、图表制作、文字润色和结果复核等辅助工作，详细使用情况见附录。

\begin{thebibliography}{99}
\setlength{\itemsep}{1.35\baselineskip}

\bibitem{qian2021energy}
钱声攀，邱奔，李哲，王少鹏．数据中心能效优化策略研究[J]．信息通信技术与政策，2021，47(4)：19--26．DOI：\url{https://doi.org/10.12267/j.issn.2096-5931.2021.04.004}．

\bibitem{li2026green}
李桃．绿色数据中心环境下大数据处理的能耗优化策略研究[J]．电脑知识与技术，2026，22(13)：56--58．DOI：\url{https://doi.org/10.14004/j.cnki.ckt.2026.0622}．

\bibitem{bergmeir2018}
Bergmeir C, Hyndman R J, Koo B. A note on the validity of cross-validation for evaluating autoregressive time series prediction[J]. \emph{Computational Statistics \& Data Analysis}, 2018, 120: 70--83. DOI: \url{https://doi.org/10.1016/j.csda.2017.11.003}.

\bibitem{hyndman2021}
Hyndman R J, Athanasopoulos G. \emph{Forecasting: Principles and Practice}[M/OL]. 3rd ed. Melbourne: OTexts, 2021. \url{https://otexts.com/fpp3/}.

\bibitem{friedman2001}
Friedman J H. Greedy function approximation: A gradient boosting machine[J]. \emph{The Annals of Statistics}, 2001, 29(5): 1189--1232. DOI: \url{https://doi.org/10.1214/aos/1013203451}.

\bibitem{hastie2009}
Hastie T, Tibshirani R, Friedman J. \emph{The Elements of Statistical Learning: Data Mining, Inference, and Prediction}[M]. 2nd ed. New York: Springer, 2009.

\bibitem{liang2020periodic}
梁毅，曾绍康，梁岩德，丁毅．一种基于周期性特征的数据中心在线负载资源预测方法[J]．计算机工程与科学，2020，42(3)：381--390．

\bibitem{wickramasuriya2019}
Wickramasuriya S L, Athanasopoulos G, Hyndman R J. Optimal forecast reconciliation for hierarchical and grouped time series through trace minimization[J]. \emph{Journal of the American Statistical Association}, 2019, 114(526): 804--819. DOI: \url{https://doi.org/10.1080/01621459.2018.1448825}.

\bibitem{vovk2005}
Vovk V, Gammerman A, Shafer G. \emph{Algorithmic Learning in a Random World}[M]. New York: Springer, 2005.

\bibitem{romano2019}
Romano Y, Patterson E, Cand\`es E J. Conformalized quantile regression[C]// \emph{Advances in Neural Information Processing Systems}. 2019, 32: 3543--3553.

\bibitem{ortools}
Perron L, Furnon V. OR-Tools: Google optimization tools[EB/OL]. Google Developers, 2026[2026-08-09]. \url{https://developers.google.com/optimization/}.

\bibitem{qian2022resource}
钱敏．云计算背景下数据中心资源调度方法研究[J]．信息技术，2022(5)：188--192．DOI：\url{https://doi.org/10.13274/j.cnki.hdzj.2022.05.032}．

\bibitem{wang2023ddpg}
王立红，张延华，孟德彬，李萌．基于 DDPG 算法的云数据中心任务节能调度研究[J]．高技术通讯，2023，33(9)：927--936．DOI：\url{https://doi.org/10.3772/j.issn.1002-0470.2023.09.004}．

\bibitem{liu2023pso}
刘陈伟，孙鉴，雷冰冰，徐涛，吴隹伟．基于改进粒子群算法的云数据中心能耗优化任务调度策略[J]．计算机科学，2023，50(7)：246--253．DOI：\url{https://doi.org/10.11896/jsjkx.220900176}．

\bibitem{sun2022aco}
孙湛冬，焦娇，李伟，李志鹏，李鹏恩．基于改进蚁群算法的电力云数据中心任务调度策略研究[J]．电力系统保护与控制，2022，50(2)：95--101．DOI：\url{https://doi.org/10.19783/j.cnki.pspc.210466}．

\bibitem{wolsey1998}
Wolsey L A. \emph{Integer Programming}[M]. New York: Wiley, 1998.

\bibitem{pinedo2016}
Pinedo M L. \emph{Scheduling: Theory, Algorithms, and Systems}[M]. 5th ed. Cham: Springer, 2016.

\bibitem{wang2017storage}
王飞，徐健，李伟，汪新浩，施啸寒．基于分布式储能系统的风储滚动优化调度方法[J]．山东大学学报（工学版），2017，47(6)：89--94．DOI：\url{https://doi.org/10.6040/j.issn.1672-3961.0.2017.208}．

\bibitem{sioshansi2009}
Sioshansi R, Denholm P, Jenkin T, Weiss J. Estimating the value of electricity storage in PJM: Arbitrage and some welfare effects[J]. \emph{Energy Economics}, 2009, 31(2): 269--277. DOI: \url{https://doi.org/10.1016/j.eneco.2008.10.005}.

\nocite{openai-codex}
\bibitem{openai-codex}
OpenAI Codex，GPT-5，OpenAI，2026-08-07 至 2026-08-09。

\end{thebibliography}
